\documentclass[11pt]{article}
\usepackage[T1]{fontenc}\usepackage{lmodern,amsmath,amssymb,amsthm,microtype}
\usepackage[margin=.85in]{geometry}\usepackage[colorlinks=true,urlcolor=blue,linkcolor=black,citecolor=black]{hyperref}
\newtheorem{theorem}{Theorem}\newtheorem{lemma}[theorem]{Lemma}
\newcommand{\C}{\mathbb C}\newcommand{\ket}[1]{|#1\rangle}\newcommand{\bra}[1]{\langle#1|}\newcommand{\Tr}{\operatorname{Tr}}
\setlength{\parindent}{0pt}\setlength{\parskip}{4pt}
\title{The Watrous two-qubit basis requires more than four collective local turns on two copies}
\author{Research manuscript}\date{20 November 2026}
\begin{document}\maketitle
\begin{abstract}\noindent Bandyopadhyay showed that two identical copies of the Watrous basis are distinguishable by collective separable measurements but not by adaptive separable measurements acting on one copy at a time. We exclude all collective LOCC protocols with four alternating local turns, in either direction. The proof constrains the first effect on its product-preserving branches to one of three coherent phases, then shows that the second laboratory's effects on those branches cannot sum to a subeffect of identity. Local operations and POVMs may act jointly across both copies. The unrestricted two-way LOCC question is not settled.\end{abstract}
\section{Model and notation}
Write $\mathcal Z=\{1,\omega,\omega^2\}$, $\omega=e^{2\pi i/3}$. The original two-qubit states are $3^{-1/2}(\ket{00}+z\ket{01}+z^2\ket{10})$, $z\in\mathcal Z$, and $\ket{11}$. Two copies carry the same unknown label. Index each lab's basis as $e_0=\ket{00},e_1=\ket{01},e_2=\ket{10},e_3=\ket{11}$, and put $s=e_1+e_2$, $h_0=e_1-e_2$. Juxtaposition denotes Alice--Bob tensor product. With a common factor $1/3$ omitted, the entangled candidates are $\Psi_z=u_0+zu_1+z^2u_2$, where
\begin{equation} u_0=e_0e_0+e_1e_2+e_2e_1,\quad u_1=e_3e_0+e_0s,\quad u_2=se_0+e_0e_3.\label{states}\end{equation}
The product candidate is $e_3e_3$. Swapping labs interchanges $u_1,u_2$ and permutes the labels. Bandyopadhyay~\cite{B}, Eq.~(6), credits the example to Watrous; Theorem~1 there gives collective SEP and Proposition~1 excludes adaptive copywise SEP.
\begin{theorem}No collective LOCC protocol of the form Alice--Bob--Alice--Bob, or the reverse, perfectly discriminates these four two-copy states.\end{theorem}
A turn is one laboratory's move, the last making the decision. This excludes protocols with fewer turns, without ruling out five or more. No extra shared entanglement is provided.

At any successful intermediate product filter $A^{1/2}\otimes B^{1/2}$ the four resulting vectors are orthogonal: a perfect POVM effect fixes its own vector and annihilates the others. Arbitrary instruments can be replaced by these square roots up to local isometries by retaining local environments. The three entangled states are pairwise orthogonal, with zeros allowed, exactly when the Gram matrix of the filtered $u_0,u_1,u_2$ is circulant, since Fourier transformation diagonalizes precisely those matrices. In particular their squared norms are equal and their cyclic inner products are equal. We take inner products linear in the second argument.

We use two elementary PPT facts. First, if positive $E,E^{\Gamma_A}$ satisfy $E(xy)=0$, then $E^{\Gamma_A}(\bar x y)=0$ by its zero expectation, hence $\langle x c|E|d y\rangle=0$ for all $c,d$ by transposing the matrix elements. There is a symmetric version for $\Gamma_B$. Second, write $u'_j$ for the filtered components and $\Psi'_z$ for the filtered states. For a perfect effect $E_z$,
\begin{equation} E_z u'_j=z^{-j}\Psi'_z/3 .\label{fourier}\end{equation}
Indeed this is Fourier inversion of the perfect discrimination conditions. Hence if $u'_j=v_j e$ and $u'_k=v_k e$ share a Bob factor, then
\begin{equation} (\bra{v_j-z^{k-j}v_k}\otimes I)\Psi'_z=0.\label{share}\end{equation}
For $E_z$ kills $(v_j-z^{k-j}v_k)e$ and maps $v_j e$ to a nonzero constant times $\Psi'_z$; apply the first fact. The identity remains true if $\Psi'_z=0$. It applies also when the two components share an Alice factor.

\section{The first effect is forced onto three phases}
Fix a first Alice effect $A\ge0$ with $t=A_{33}>0$; such outcomes are needed by completeness. Write
\[
 a=A^{1/2}e_0,\quad x=\|a\|^2,\quad V_i=A^{1/2}e_i\ (i=1,2),\quad
 n=A^{1/2}e_3,\quad p_A=V_1+V_2 .
\]
Perfect orthogonality to the product at $(A,I)$ gives $A_{30}=0$. The equal norms and equal cyclic inner products at this branch, directly from~\eqref{states}, give
\begin{equation}
 \|p_A\|^2=x+t=\|V_1\|^2+\|V_2\|^2,\qquad
 \langle p_A,a\rangle=\langle n,p_A\rangle.\label{rootg}
\end{equation}
In particular $\operatorname{Re}\langle V_1,V_2\rangle=0$. We must have $x>0$: otherwise $u'_0=V_1e_2+V_2e_1$, $u'_1=ne_0$, $u'_2=p_Ae_0$, and projecting~\eqref{share} on Bob coordinates 1 and 2 gives $\langle n-zp_A,V_j\rangle=0$ for all $j,z$. Thus $p_A\perp V_1,V_2$ though $\|p_A\|^2=t$. Such a first outcome cannot be completed even by PPT.

For Bob's next effect $B\ge0$ write $w_i=B^{1/2}e_i$, $p_B=w_1+w_2$, and consider outcomes with $B_{33}>0$. Orthogonality to the product for $u'_1=n w_0+a p_B$ implies $w_0\perp w_3$ since $a\perp n$. Equality of $\|u'_1\|,\|u'_2\|$ and~\eqref{rootg} gives
\begin{equation} \|p_B\|^2=B_{00}+B_{33}.\label{corner}\end{equation}
All these equalities are linear in $B$, so they also hold upon summing Bob effects.

\begin{lemma}There is a Bob outcome with both $B_{33}>0$ and $B_{00}>0$.\end{lemma}
\begin{proof}Otherwise pool all Bob effects with $B_{33}>0$ into $B_*$. Since all other effects kill $e_3$, $0\le B_*\le I$, $B_*e_3=e_3$, and $B_*e_0=0$. A successful continuation is possible from the pooled outcome: Bob can first apply $B_*^{1/2}$ and then refine to the original effects using the inverse on its support. Thus it too has a perfect PPT continuation. In this proof use $w_i=B_*^{1/2}e_i$. The vector $p_B$ has squared norm one by~\eqref{corner}, and $w_3=e_3\perp w_1,w_2$. Here $u'_1=a p_B$, $u'_2=a w_3$. Applying the Bob version of~\eqref{share} to these two and varying $z$ yields
\[
        \langle p_B,w_2\rangle V_1+\langle p_B,w_1\rangle V_2=0.
\]
Indeed the $z$-dependent contribution from $z a p_B+z^2 a w_3$ cancels since the Bob vectors are orthonormal. The two coefficients sum to one. Thus $V_1=r g,V_2=q g$ for some unit $g$ and $r,q$ not both zero. By~\eqref{rootg}, $\operatorname{Re}(\bar r q)=0$, so $r-q\ne0$. The two coefficients must be $-q/(r-q)$ and $r/(r-q)$. Their squared moduli sum to one. Therefore $\|B_*s\|^2=1=\langle s,B_*s\rangle$ (the corner rows are zero on $s$). Since $B_*\le I$, $B_*s$ and $(I-B_*)s$ are orthogonal eigenvectors of eigenvalues 1 and 0 respectively. Both have norm one and span the middle two-dimensional block. Hence $B_*$ is rank one there, $w_1,w_2$ are collinear with $p_B$, and the displayed equation forces $u'_0=V_1w_2+V_2w_1=0$, while $u'_1$ is nonzero. This contradicts equal component norms.\end{proof}

For such an outcome with both corners positive, the span of $w_0,w_1,w_2$ has dimension at least two. Otherwise $w_1=r w_0,w_2=q w_0$ and $u'_0=(a+q V_1+r V_2)w_0$, $u'_1=(n+(r+q)a)w_0$ share their Bob factor. The component of $u'_2$ along $w_3\perp w_0$ is $a w_3$. By~\eqref{share} for $j=0,k=1$, variation of $z$ forces $a\perp n+(r+q)a$. Thus $r+q=0$, contradicting~\eqref{corner}.

Alice acts once more and Bob makes the final decision. Alice's measurement may be spectrally refined to rank-one effects, giving total accumulated bras with coefficients $r_i$ satisfying $\sum_r\overline{r_i}r_j=A_{ij}$. Indeed the conditional Bob operators after tracing out Alice are positive linear functions of each Alice effect, independent of the instrument; an error-free final Bob decoder for their sum has zero wrong-label probability on every positive summand. The rank-one effects here are in the original coordinates, after pulling the refinement back through $A^{1/2}$.

If $r_3\ne0$ then the product survives, forcing $r_0=0$ by the $u'_2$ component and $w_3\perp w_0$. The other three outputs are
\begin{equation}r_2w_1+r_1w_2+\bigl(z r_3+z^2(r_1+r_2)\bigr)w_0.\label{last}\end{equation}
With $w_0\ne0$, their span has dimension at most two. One must vanish; this makes them collinear, so two must vanish. The displayed Fourier coefficients are not all equal. If $z_0$ is the surviving label, coincidences at the other two give $r_1+r_2=z_0^2 r_3$, with common coefficient $-z_0 r_3$. Thus
\[
   -z_0 r_3 e_0+r_2 e_1+r_1 e_2\in\ker B.
\]
The kernel on $\operatorname{span}(e_0,e_1,e_2)$ has dimension at most one, as just proved. Its ray therefore fixes $z_0$ (the negative ratio of the first coordinate to the sum of the next two is $z_0^2$), and then fixes the full ray of $(r_0,r_1,r_2,r_3)$. Summing against $\bar r_3$ gives
\begin{equation} A_{3s}/t=\langle n,p_A\rangle/t\in\mathcal Z .\label{phase}\end{equation}

This holds for every first effect with $t>0$. Project $p_A$ onto the orthogonal vectors $a,n$. By~\eqref{rootg} and~\eqref{phase}, $x+t=\|p_A\|^2\ge t^2/x+t$, so $x\ge t$. The first measurement has $\sum x=\sum t=1$, forcing $x=t$ throughout and equality in this projection. Choose any first outcome with $t>0$ and rescale $A$ by $1/t$, which does not affect conditional distinguishability. By the local symmetry that multiplies a one-copy Alice $\ket1$ by $\omega^2$ and a Bob $\ket1$ by $\omega$, the phase in~\eqref{phase} can be made one, relabeling $z$ cyclically. It follows that
\begin{equation}
  \langle a,n\rangle=0,\quad \|a\|=\|n\|=1,\qquad
  p_A=a+n,\qquad d:=V_1-V_2,\quad \|d\|^2=2.\label{sat}
\end{equation}
The final identity follows by $\operatorname{Re}\langle V_1,V_2\rangle=0$.

\section{Pooling the surviving Bob effects}
Set $\alpha=\langle a,d\rangle$, $\beta=\langle n,d\rangle$, so $|\alpha|^2+|\beta|^2\le2$. Consider \emph{every} Bob outcome with $B_{33}>0$, including $B_{00}=0$. For each final Alice bra with $r_3\ne0$, we still have $r_0=0$ since $w_0\perp w_3$, and $r_1+r_2=r_3$ by~\eqref{sat} (the bras are contractions with the filtered Alice vectors). In~\eqref{last} the outputs at $z=\omega,\omega^2$ are equal and orthogonal, hence zero, whether or not $w_0$ vanishes. Multiply this identity by $\bar r_3$ and sum over the final Alice outcomes. It gives
\[
  B^{1/2}\bigl(-e_0+\langle n,V_2\rangle e_1+\langle n,V_1\rangle e_2\bigr)=0 .
\]
The vector in the kernel is the same for all these Bob outcomes. Now let $B$ denote their sum, so $0\le B\le I$, $Be_3=e_3$, and $B$ annihilates the same vector. Put $b=B^{1/2}e_0$, $p=B^{1/2}s$, $h=B^{1/2}h_0$, $m=e_3$. Then
\begin{equation} p=2b+\beta h,\quad m\perp b,p,h,\quad
 X:=\|b\|^2,\quad Z:=\|h\|^2\le2,\quad Y:=\langle b,h\rangle,\quad \|p\|^2=1+X.\label{geom}\end{equation}
The filtered Fourier components for this pooled effect are
\[
 U_0=ab+\tfrac12[(a+n)p-dh],\qquad U_1=nb+ap,\qquad U_2=(a+n)b+am .
\]
Their Gram $G_{ij}=\langle U_i,U_j\rangle$ is circulant by linearity in $B$.

\section{Scalar contradiction}
For transparency put $v=|\beta|^2$, $c+i\sigma=\bar\alpha\beta$, $L=2X+\beta Y$, $H=2\bar Y+\beta Z$, and $N=3X+4\operatorname{Re}(\beta Y)+vZ-1$. Eq.~\eqref{geom} gives $N=0$. Direct expansion of the displayed $U_i$ gives
\begin{align*}
 G_{20}-G_{12}&=L-\bar L-(\alpha+\beta)Y/2,\\
 2(G_{01}-G_{12})-N&=1+X+2(\beta Y-\overline{\beta Y})-\bar\alpha H,\\
 2(G_{00}-G_{11})+N&=3X-4\operatorname{Re}(\alpha Y)+(1-c-v)Z-1 .
\end{align*}
Every left side is zero. In particular $Z>0$, since otherwise $Y=H=0$ and the middle equation gives $1+X=0$. Taking the real and imaginary parts of the first equation gives
\[
  \operatorname{Re}[(\alpha+\beta)Y]=0,\qquad
  \operatorname{Im}[(\alpha-3\beta)Y]=0.
\]
Subtract $N=0$ from the final equation, then use the real relation to obtain $c+2v=1$. Thus $\beta\ne0$ (or both $c$ and $v$ would vanish). Put $k=\beta Y/Z$. Separating the middle equation and using the two relations gives
\begin{align}
 &\operatorname{Re}[(\alpha/\beta+1)k]=0,\quad
   \operatorname{Im}[(\alpha/\beta-3)k]=0,\quad
   \operatorname{Im}k=\sigma/10,\label{red1}\\
 &1+X=(c-2\operatorname{Re}k)Z,\qquad
   1=3X+(v+4\operatorname{Re}k)Z .\label{red2}
\end{align}
If $\sigma\ne0$, then $\alpha/\beta=(c-i\sigma)/v$. The imaginary relation in~\eqref{red1} gives $10\operatorname{Re}k=c-3v=1-5v$ and the real relation gives $\sigma^2=(1-v)(5v-1)>0$, so $v>1/5$. Substitution into~\eqref{red2} yields $Z=10/(7-10v)$; $0<Z\le2$ forces $v\le1/5$, a contradiction.

If $\sigma=0$, $k$ is real. When $c\ne-v$, Eq.~\eqref{red1} gives $k=0$, so $Z=4/(3-5v)$. Since $Z\le2$ we have $v\le1/5$. But $|\alpha|^2+v\le2$ gives $(1-2v)^2/v+v\le2$, or $(5v-1)(v-1)\le0$, so $v=1/5$ and $Z=2$. Equality implies $Bh_0=h_0$ since $B\le I$; hence $\langle h,p\rangle=\langle h_0,s\rangle=0$. On the other hand $Y=0$ and $p=2b+\beta h$ give $\langle h,p\rangle=\beta Z\ne0$.

Finally if $\sigma=0$ and $c=-v$, then $v=1$ and $\alpha=-\beta$. The first identity in~\eqref{red2} gives $1+X+Z=-2\operatorname{Re}(\beta Y)\le2\sqrt{XZ}$, another contradiction. This exhausts the cases, proving the theorem. No numerical SDP premise is used.

For measurable outcomes all effect equalities can be read for densities relative to trace measure. Perfect conditional action holds almost everywhere by nonnegativity of errors. Pooling is integration of positive effects, kernel constraints transfer under that integration, and completeness requires positive total weight on the indicated surviving sets.

\begin{thebibliography}{9}
\bibitem{B} S.~Bandyopadhyay, \emph{Perfect discrimination of an orthogonal $2\otimes2$ basis may require collective separable measurements on two copies}, \href{https://arxiv.org/abs/2610.16138v2}{arXiv:2610.16138v2} (7 November 2026), Eqs.~(4)--(6), Proposition~1, Theorem~1 and Section~7.
\end{thebibliography}
\end{document}
